% =====================================================================
% Data Centers Orbitais de IA: Estrutura de LCA e Limiares Parametricos
% Compilar com: pdflatex (2x) ou xelatex
% Figuras: python odc_lca_figures.py (gera os PDFs na mesma pasta)
% =====================================================================
\documentclass[11pt,a4paper]{article}

\usepackage[utf8]{inputenc}
\usepackage[T1]{fontenc}
\usepackage[brazilian]{babel}
\usepackage{lmodern}
\usepackage{amsmath,amssymb}
\usepackage{booktabs}
\usepackage{longtable}
\usepackage{array}
\usepackage{graphicx}
\usepackage{microtype}
\usepackage[margin=2.5cm]{geometry}
\usepackage{setspace}
\usepackage{caption}
\usepackage[hidelinks,breaklinks]{hyperref}

\setstretch{1.08}
\captionsetup{font=small,labelfont=bf}
\renewcommand{\arraystretch},{1.15}

\newcommand{\CO},{\ensuremath{\mathrm{CO_2}}}
\newcommand{\COeq},{\ensuremath{\mathrm{CO_2}\text{-eq}}}

% inclui a figura se o arquivo existir; caso contrario, mostra um aviso
\newcommand{\figinc}[2]{%
  \IfFileExists{#1},{\includegraphics[width=#2]{#1}},{%
    \fbox{\parbox{0.8\textwidth},{\centering\small
    [Figura ausente: execute \texttt{odc\_lca\_figures.py}]}}}}

\title{\textbf{Data Centers Orbitais de Inteligência Artificial:\\
Estrutura de Avaliação de Ciclo de Vida e Limiares Paramétricos\\
de Equilíbrio de Carbono e Água}}

\author{
  Dafnhê Thor Castro de Martins Barros\thanks{Mestrando em Engenharia de Controle e Automação, Instituto Federal do Espírito Santo (IFES), Brasil. E-mail: dafnhe.barros@aluno.ifes.edu.br}
  \and
  Flávio Barcelos Braz da Silva\thanks{Doutor em Engenharia Elétrica, Professor Associado do PROPECAUT/IFES, Brasil. E-mail: flavio.silva@ifes.edu.br}
}
\date{Outubro de 2026}

\begin{document}
\maketitle

% ---------------------------------------------------------------------
\begin{abstract}
\noindent
O crescimento da inteligência artificial (IA) intensifica a demanda elétrica e
hídrica de \emph{data centers} terrestres e estimula a investigação de
\emph{data centers} orbitais (ODCs). A irradiância solar extraterrestre e a
ausência de torres evaporativas conferem vantagens operacionais, mas não
demonstram superioridade ambiental no ciclo de vida. Este trabalho propõe uma
estrutura de Avaliação de Ciclo de Vida (LCA) alinhada à ISO 14040/14044 e a
combina com um modelo paramétrico de primeira ordem de um nó orbital de 1~MW de
potência de TI em órbita heliossíncrona (SSO) a 700~km, complementar a modelos
existentes de nó de 1~kW. O modelo trata geometria de eclipse, potência,
armazenamento, radiadores, momento de atitude por gradiente de gravidade,
arrasto e confiabilidade. Os resultados são cenários condicionais, não
medições. A unidade funcional provisória (1~MW de TI por dez anos) corresponde
a 87,6~GWh de energia de TI e a 105,12--131,4~GWh de eletricidade total com PUE
de 1,2--1,5. O consumo direto de água terrestre é 28,0--58,7 milhões de litros
(WUE de 0,32--0,67~L/kWh), 87,6--175,2 milhões (WUE de 1--2~L/kWh, cenário dos
autores) e 438 milhões (5~L/kWh). A água da geração elétrica é de ordem
comparável à direta para matriz eólica e supera-a em mais de duas ordens de
grandeza para matriz hidrelétrica. A massa de lançamento estimada é 52,4~t
(órbita crepuscular, LTAN 06:00) e 68,3~t (LTAN $\approx$ 12:00), com
intervalos de 90\% de 57,2--91,3~t e 73,0--111,1~t pelo método de Monte Carlo.
A massa areal do radiador, a das mantas fotovoltaicas e a massa específica da
eletrônica dominam a incerteza, e o resultado depende da distribuição adotada
para a massa areal do radiador. Sob a premissa adotada, a confiabilidade exige
sobredimensionamento de 11\% a 76\% das unidades de computação. A
Tabela~\ref{tab:teto} fornece tetos de emissões incorporadas diferenciais. Um
piso físico de $\CO$ de combustão de lançamento (6,1--15,6~kg $\CO$/kg de
carga) não esgota esses tetos; contudo, com as intensidades de lançamento e
reentrada de Ohs \emph{et al.} (34,4--51,7~kg $\COeq$/kg), esses dois termos
ocupam de 3\% a 7,5\% do teto em rede de 450~g/kWh e de 27\% a 67\% em rede de
50~g/kWh, podendo superá-lo no P95 do caso B. Conclui-se que a vantagem
ambiental de um ODC permanece uma hipótese condicional.

\medskip
\noindent\textbf{Palavras-chave:} avaliação de ciclo de vida; \emph{data
centers} orbitais; inteligência artificial; lançadores espaciais; pegada de
carbono; consumo de água.
\end{abstract}

\begin{center}
\textbf{Abstract}
\end{center}
\begin{quote}\small
\noindent
The growth of artificial intelligence (AI) increases the electricity and
water demand of terrestrial data centers and motivates the investigation of
orbital data centers (ODCs). Extraterrestrial solar irradiance and the absence
of evaporative cooling towers give operational advantages, but do not by
themselves demonstrate life-cycle environmental superiority. This work
proposes a Life Cycle Assessment (LCA) framework aligned with ISO
14040/14044 and combines it with a first-order parametric model of a 1~MW (IT
power) orbital node in a 700~km Sun-synchronous orbit (SSO), complementary to
existing 1~kW node models. The model covers eclipse geometry, power
generation, energy storage, radiators, gravity-gradient momentum, drag, and
reliability. Results are conditional scenarios, not measurements. The
provisional functional unit (1~MW of IT power for ten years) corresponds to
87.6~GWh of IT energy and to 105.12--131.4~GWh of total electricity for a PUE
of 1.2--1.5. Direct terrestrial water consumption is 28.0--58.7 million litres
(WUE of 0.32--0.67~L/kWh), 87.6--175.2 million (WUE of 1--2~L/kWh, authors'
scenario) and 438 million (5~L/kWh). Water use in electricity generation is
comparable to direct use for a wind-based grid and exceeds it by more than two
orders of magnitude for a hydropower-based grid. The estimated launch mass is
52.4~t (dawn--dusk orbit, LTAN 06:00) and 68.3~t (LTAN $\approx$ 12:00), with
90\% intervals of 57.2--91.3~t and 73.0--111.1~t from Monte Carlo propagation.
Radiator areal mass, solar-array areal mass and electronics specific mass
dominate the uncertainty, and the result depends on the distribution assumed
for the radiator areal mass. Under the adopted premise, reliability requires
oversizing the compute units by 11\% to 76\%. Table~\ref{tab:teto} provides
ceilings on differential embodied emissions. A physical floor on launch
combustion $\CO$ (6.1--15.6~kg $\CO$/kg of payload) does not exhaust these
ceilings; however, applying the launch and re-entry intensities of Ohs
\emph{et al.} (34.4--51.7~kg $\COeq$/kg), these two terms take 3\% to 7.5\% of
the ceiling on a 450~g/kWh grid and 27\% to 67\% on a 50~g/kWh grid, and could
exceed it at the P95 of case B. We conclude that the environmental advantage of
an ODC remains a conditional hypothesis.

\medskip
\noindent\textbf{Keywords:} life cycle assessment; orbital data centers;
artificial intelligence; space launch; carbon footprint; water use.
\end{quote}

\bigskip
% ---------------------------------------------------------------------
\section{Introdução}

A IA de grande escala transforma a infraestrutura computacional em fator
relevante de planejamento energético. A Agência Internacional de Energia
estima que os \emph{data centers} consumiram cerca de 415~TWh em 2024
($\approx$1,5\% da eletricidade mundial) e projeta cerca de 945~TWh em 2030 no
cenário de referência~\cite{r5}. Essas projeções são cenários. Os relatórios do
LBNL mostram a incerteza substancial da trajetória da demanda~\cite{r6,r7}.

A concentração regional de instalações pressiona redes e recursos hídricos. O
Öko-Institut~\cite{r8} relata participações de 26\% (Virgínia) e 21\% (Irlanda)
no consumo elétrico regional em 2023, e cerca de 80\% em Dublin, com ressalvas
sobre fontes distintas (a IEA estima 17\% para a Irlanda).

Propostas recentes cogitam levar \emph{data centers} de IA ao espaço. Em 8 de
julho de 2026, DarkSky International, Environment America e PEER peticionaram à
FCC a elaboração de um \emph{Programmatic Environmental Impact Statement} (PEIS)
sob o NEPA para pedidos de licenciamento de \emph{data centers} orbitais, e a
suspensão do licenciamento até sua conclusão~\cite{r50}; a petição foi
noticiada pela imprensa~\cite{r24}, e o custo ambiental da proposta também é
discutido na imprensa especializada~\cite{r23}. Um \emph{white paper} setorial
discute o racional, os casos de uso, a arquitetura de referência e os custos de
\emph{data centers} espaciais~\cite{r20}.

O lançamento não é neutro: a combustão injeta produtos em regiões sensíveis da
atmosfera, e a reentrada produz material ablacionado~\cite{r13,r14,r15,r16}.
Uma comparação válida precisa incluir fabricação, lançamento, operação,
reposição e fim de vida. Al~Ahmad \emph{et al.}~\cite{r17} analisam riscos e
confiabilidade no ciclo de vida de infraestruturas em nuvem orbitais, sem
substituir um inventário quantitativo.

\paragraph{Trabalhos mais próximos.} Os precedentes diretos são de Ohs
\emph{et al.} O ESpaS~\cite{r18} estima a intensidade de carbono de CPU,
memória e rede em órbita e em solo e conclui que sistemas em órbita têm custo
de carbono significativamente maior, principalmente por emissões incorporadas
de lançamento e reentrada. Sua extensão ESpaS-ODC~\cite{r19} acrescenta, para
um nó de 1~kW com uma GPU H100 a 510~km, um radiador calibrado com dados da
ISS, a degradação do arranjo solar com margem de eclipse e sobressalentes em
\emph{standby} frio, e reporta o carbono por GPU-hora. Na fronteira adotada
pelos autores, o nó orbital alcança a paridade com um \emph{data center}
terrestre de média global em até cerca de dois anos de missão, mas a paridade
com um \emph{data center} ``verde'' exige missões de seis a nove anos com
sobressalentes~\cite{r19}. O modelo de~\cite{r19} deixa de fora o barramento
(estrutura, ADCS, aviônica) e considera, na fronteira terrestre, apenas o
carbono incorporado do \emph{chip}. Ficam em aberto: (a) o dimensionamento de
um nó de escala de megawatt com barramento, ADCS e eclipse dependente da
órbita; (b) o balanço hídrico; (c) a propagação da incerteza dos parâmetros de
massa; e (d) a relação entre um teto de emissões incorporadas e a intensidade
de carbono da rede.

\paragraph{Contribuições.} (i)~Estrutura de LCA comparativa com fronteiras,
equações de balanço e critérios de alocação. (ii)~Modelo paramétrico
reprodutível de um nó de 1~MW em SSO que inclui o que~\cite{r19} deixa fora do
nó de 1~kW (eclipse por LTAN, momento de atitude por gradiente de gravidade,
arrasto, estrutura e margem de massa), verificado em tempo discreto e comparado
quantitativamente com os parâmetros de~\cite{r19} (Seção~\ref{sec:atm} e
Tabela~\ref{tab:cal19}). O modelo é complementar ao ESpaS-ODC: não calcula
carbono incorporado nem carbono por GPU-hora. (iii)~Tetos de emissões
incorporadas diferenciais e balanço hídrico direto e indireto, com
sensibilidade e propagação de incerteza. (iv)~Lista de requisitos de dados e de
GNC, potência e estimação para transformar o \emph{framework} em LCA
verificada.

\paragraph{Escopo.} O trabalho não é uma LCA verificada de um ODC comercial,
nem propõe uma lei de controle. As afirmações de engenharia limitam-se ao que o
modelo paramétrico sustenta.

% ---------------------------------------------------------------------
\section{Metodologia de LCA}

\subsection{Referencial normativo}\label{sec:norma}

A avaliação segue os princípios e requisitos da ISO 14040:2006~\cite{r1} e da
ISO 14044:2006, com as emendas A1:2017 e A2:2020~\cite{r2}. As categorias de
impacto devem ser coerentes com objetivo e escopo e abrangentes para uma
comparação pública. Alegações comparativas divulgadas ao público exigem revisão
crítica por painel de partes interessadas (ISO 14044, Seção~6), o que não foi
realizado aqui. Por isso, nenhum resultado deste artigo deve ser divulgado como
alegação de superioridade ambiental. A estrutura segue a prática de LCA
descrita em~\cite{r37} e trata a água de forma regionalizada
conforme~\cite{r38,r39}.

\subsection{Objetivo e unidade funcional}

O objetivo é delimitar as condições sob as quais um sistema orbital pode
igualar ou superar um sistema terrestre na entrega contínua de computação
durante dez anos.

\textbf{Unidade funcional provisória (UF):} operação contínua de 1~MW de
potência elétrica de TI por dez anos. Ela normaliza os cálculos de energia e
água e não demonstra equivalência de computação útil. A equivalência exige uma
carga de trabalho de referência: taxa de inferência ou treino, modelo,
precisão numérica, utilização, latência, disponibilidade e volume de dados
transferido. A métrica de pegada funcional é
\begin{equation}
  I_{GWP}=\frac{G_{lifecycle}},{Q},
  \label{eq:igwp}
\end{equation}
em que $Q$ é a quantidade de trabalho útil entregue.

\subsection{Fronteiras do sistema}
\begin{enumerate}
  \item Matérias-primas e processamento.
  \item Semicondutores, memória e encapsulamento.
  \item Painéis, baterias, radiadores, estrutura e eletrônica.
  \item Integração, ensaios, transporte e preparação do lançamento.
  \item Propelentes, operação do lançador e injeção orbital.
  \item Operação, armazenamento e gerenciamento térmico.
  \item Redundância, degradação, substituição e perda de hardware.
  \item Desativação, reentrada ou órbita cemitério.
  \item Tratamento de resíduos.
\end{enumerate}
O tratamento de fim de vida deve registrar separadamente reentrada controlada,
reentrada não controlada, transferência para órbita cemitério e permanência sem
manobra. Atribuir impacto zero a qualquer uma delas exige justificativa.

\subsection{Categorias de impacto}

Potencial de aquecimento global (GWP, kg~$\COeq$); energia primária (MJ); água
(retirada, consumo, escassez); carbono negro e poluentes de lançamento; ozônio;
partículas e ablação de reentrada; matérias-primas críticas e resíduos.

O GWP convencional não representa o efeito de emissões em altitude. Forçante
radiativa de carbono negro, química do ozônio e deposição de alumina exigem
modelos atmosféricos específicos e não devem ser somados a kg~$\COeq$ sem
método de caracterização documentado.

% ---------------------------------------------------------------------
\section{Modelo paramétrico do nó orbital de 1~MW}

Os parâmetros de entrada são hipóteses de projeto conceitual
(Tabela~\ref{tab:premissas}) e não resultados de \emph{datasheets}. As normas
AIAA S-120A~\cite{r3} e ECSS-E-ST-10-02C~\cite{r4} aparecem como referências de
rastreabilidade de engenharia. Não validam os parâmetros de massa.

\subsection{Geometria orbital e eclipse}

Para $h=700$~km, $r=R_E+h=7078{,}1$~km e $\mu=398600{,}4$~km$^3$/s$^2$:
\begin{equation}
  T=2\pi\sqrt{\frac{r^{3}},{\mu}}=98{,}77~\text{min},\qquad
  \cos i_{SSO}=-\frac{\dot\Omega_{SSO}},{\tfrac32\, n\, J_2 (R_E/r)^2}
  \;\Rightarrow\; i=98{,}19^\circ .
  \label{eq:orb}
\end{equation}
Com sombra cilíndrica~\cite{r25,r26}, a duração do eclipse para ângulo beta
$|\beta|<\beta^{*}=\arcsin(R_E/r)=64{,}3^\circ$ é
\begin{equation}
  T_{ecl}=\frac{T},{\pi}\arccos\!\left(\frac{\sqrt{h^{2}+2R_Eh}},{r\cos\beta}\right).
  \label{eq:ecl}
\end{equation}
Para $|\beta|\ge\beta^{*}$ não há eclipse. Dois casos delimitam o problema:
\begin{itemize}
  \item \textbf{Caso A (órbita crepuscular, LTAN 06:00).} O $|\beta|$ anual
  varia de $58{,}4^\circ$ a $89{,}9^\circ$. Há eclipse em 83 dias por ano, com
  máximo de \textbf{18,8~min}; a temporada de eclipse ocorre aproximadamente de
  9 de novembro a 30 de janeiro. Em dez anos, cerca de 12.100 das 53.249 órbitas
  sofrem eclipse ($\sim$23\%).
  \item \textbf{Caso B (LTAN $\approx$ 12:00, $\beta\approx0^\circ$).} O eclipse
  é de \textbf{35,3~min} em todas as órbitas (53.249 ciclos em dez anos).
\end{itemize}
\emph{Hipóteses:} sombra cilíndrica, sem penumbra; Sol em movimento médio;
obliquidade de $23{,}44^\circ$. A LTAN de 18:00 produz o mesmo número de dias e
os mesmos valores, com a temporada deslocada em meio ano. As
Figuras~\ref{fig:beta} e~\ref{fig:eclvsbeta} mostram a evolução anual de
$|\beta|$ e a duração do eclipse.

\begin{figure}[ht]
\centering
\figinc{fig_beta_eclipse.pdf},{0.85\textwidth}
\caption{Ângulo beta anual (em módulo) e duração do eclipse por órbita no caso
A (LTAN 06:00), SSO a 700~km. A região sombreada indica os dias com
$|\beta|<\beta^{*}$; como a temporada de eclipse atravessa a virada do ano, ela
aparece nas duas extremidades do eixo. A linha pontilhada inferior marca o
eclipse constante do caso B. Sol em movimento médio e sombra cilíndrica.}
\label{fig:beta}
\end{figure}

\begin{figure}[ht]
\centering
\figinc{fig_eclipse_vs_beta.pdf},{0.7\textwidth}
\caption{Duração do eclipse em função de $|\beta|$ para $h=700$~km. Os pontos
indicam o máximo anual do caso A ($|\beta|=58{,}4^\circ$; 18,8~min) e o caso B
($\beta\approx0^\circ$; 35,3~min).}
\label{fig:eclvsbeta}
\end{figure}

\subsection{Potência e armazenamento}\label{sec:potencia}

Com $P_{load}=1{,}05$~MW (1~MW de TI mais 5\% de aviônica e térmica; P1),
$\eta_{conv}=0{,}97$ e $\eta_{rt}=0{,}92$ (P2):
\begin{align}
  E_{ecl}&=P_{load}\,T_{ecl}, \quad
  E_{bat}=\frac{E_{ecl}},{DoD},\quad
  M_{bat}=\frac{E_{bat}},{e_b},\quad DoD=0{,}35,\; e_b=165~\text{Wh/kg},
  \label{eq:bat}\\
  P_{gen}&=\frac{P_{load}},{\eta_{conv}}
   +\frac{E_{ecl}},{\eta_{conv}\,\eta_{rt}\,T_{sun}},
  \label{eq:pgen}\\
  A_{PV}&=\frac{P_{gen}},{\eta_{cell}\,k_{arr}\,G_{min}},\qquad
  M_{PV}=\rho_{PV}\,A_{PV}.
  \label{eq:apv}
\end{align}
Em~\eqref{eq:apv}, $G_{min}=G_0/(1+e)^2=1316{,}7$~W/m$^2$ (afélio;
$G_0=1361$~W/m$^2$, $e=0{,}0167$~\cite{r9}), $\eta_{cell}=0{,}295$ (BOL, P4) e
$k_{arr}=0{,}80$ (empacotamento, temperatura, degradação de dez anos e
apontamento; hipótese dos autores). A densidade areal é $\rho_{PV}=3{,}2$~kg/m$^2$
(P5; faixa 2,5--4,5), o que corresponde a cerca de 97~W/kg de potência
específica.

\paragraph{Verificação em tempo discreto.} O balanço foi integrado com
$\Delta t=10$~s ao longo de cinco órbitas. O estado de carga oscila entre 0,65 e
1,0 do máximo e retorna a 0,651--0,652 ao final de cada órbita, em ambos os
casos, com SOC inicial de 0,65 (DoD de 35\%). Isso confirma o fechamento do
balanço energético, com erro de quantização da fronteira de eclipse
$\le 0{,}9\%$ ($\Delta t/T_{ecl}$ para 18,8~min). As taxas-C resultantes são:
carga de 0,26--0,33~C; descarga de 1,12~C (A) e 0,60~C (B).

\subsection{Térmica}\label{sec:termica}

A emissão de corpo cinzento é
$q_{ideal}=\varepsilon\,\eta_{fin}\,\sigma T^{4}=543$~W/m$^2$ por face
($\varepsilon=0{,}88$, $\eta_{fin}=0{,}94$, $T=328$~K). Para radiador com a
normal perpendicular ao nadir, o fator de forma da Terra por face é
$F=\tfrac{2},{\pi}\left(\rho/2-\sin 2\rho/4\right)=0{,}233$, com
$\rho=\arcsin(R_E/r)$. Sob IR planetário médio de 237~W/m$^2$ e albedo (P20), o
fluxo líquido é $\approx 485$~W/m$^2$ por face. O fluxo de projeto
$q=360$~W/m$^2$ (adotado sem fonte primária) equivale a um fator de margem de
0,74. A área e a massa são
\begin{equation}
  \bar Q=P_{gen}\frac{T_{sun}},{T},\qquad
  A_{rad}=\frac{\bar Q},{2q},\qquad
  M_{rad}=\rho_{rad}A_{rad},\quad \rho_{rad}=6~\text{kg/m}^2 .
  \label{eq:rad}
\end{equation}
A massa areal $\rho_{rad}$ é o parâmetro de maior incerteza do modelo (P6).
Adota-se a faixa de 4--22~kg/m$^2$, com moda em 6~kg/m$^2$; o limite superior
cobre o equivalente de 21,8~kg/m$^2$ obtido dos parâmetros de~\cite{r19}
(Seção~\ref{sec:atm}). No caso base, o radiador responde por 17,7\% (A) e
13,8\% (B) da massa seca.

\emph{Limitação:} para o caso B mantém-se a mesma orientação de radiador, o que
subestima a carga térmica ambiental. A mesma $\rho_{rad}$ é aplicada à área em
planta.

\subsection{Atitude e perturbações}

Os arranjos assumidos têm duas asas simétricas de 30~m de largura e cubo central
de 10~m (P11). O momento de inércia sobre o eixo transversal é
$I_{xx}\approx 3{,}0\times10^{7}$~kg\,m$^2$ (A) e $6{,}0\times10^{7}$~kg\,m$^2$
(B). O torque máximo de gradiente de gravidade e o momento cíclico necessário
são
\begin{equation}
  T_{gg}=\tfrac32\,\frac{\mu},{r^{3}}\,\Delta I,\qquad
  h_{cyc}=\frac{T_{gg}},{2n},\qquad
  M_{ADCS}=\frac{N_{red}\,h_{cyc}},{h_{spec}},
  \label{eq:adcs}
\end{equation}
com $\Delta I\approx I_{xx}$ (limite superior), $n=2\pi/T$, $N_{red}=2$ e
$h_{spec}=17{,}5$~N\,m\,s/kg (P10). Resultados: $T_{gg}=51$~N\,m (A) e 102~N\,m
(B); $h_{cyc}=2{,}4\times10^{4}$~N\,m\,s (A) e $4{,}8\times10^{4}$~N\,m\,s (B);
$M_{ADCS}=2{,}7$~t (A) e 5,5~t (B).

As demais perturbações são menores em ordem de grandeza. A força de pressão de
radiação solar é de 0,033~N (coeficiente de 1,3; área do arranjo). O torque
aerodinâmico é de $\le 0{,}17$~N\,m (braço de 5~m, hipótese). O torque magnético
para um dipolo residual de $10^{3}$~A\,m$^2$ (hipótese) é de 0,035~N\,m. A
dominância do gradiente de gravidade por duas a três ordens de grandeza
justifica o dimensionamento acima.

\emph{Limitações:} o modelo não inclui modos flexíveis dos arranjos (típicos de
baixa frequência, a determinar por análise modal) nem a gestão de momento por
equilíbrio de torque (TEA), que pode reduzir $h_{cyc}$. Uma simulação
6-DOF~\cite{r27,r31} é requisito para validar \eqref{eq:adcs}.

\paragraph{Arrasto.} $F_d=\tfrac12\rho v^{2}C_dA$ com $\rho=10^{-13}$~kg/m$^3$,
$C_d=2{,}2$, $A=A_{PV}$ e $v=7{,}5$~km/s (P13), o que resulta em
$F_d\approx0{,}03$~N. Ao longo de dez anos, isso exige
$\Delta v\approx160$--$163$~m/s. Com propulsão elétrica ($I_{sp}=1500$~s,
reserva de 10\%), a massa de propelente é $m_p\approx0{,}64$~t (A) e 0,81~t (B).
A área do radiador não entra no cálculo do arrasto.

\subsection{Massa e lançamentos}
\begin{equation}
  M_{dry}=(1+\mu_{mga})\big(M_{PV}+M_{rad}+M_{bat}+M_{IT}+M_{st}+M_{ADCS}+M_{hw}\big),
  \qquad M_{launch}=M_{dry}+m_p,
  \label{eq:mass}
\end{equation}
com $M_{IT}=6{,}0$~kg/kW $\times\,1000$~kW (P7),
$M_{st}=0{,}15\,(M_{PV}+M_{rad}+M_{bat}+M_{IT})$ (P8), $M_{hw}=300$~kg e
$\mu_{mga}=0{,}20$. A margem de 20\% é adotada como hipótese de projeto
conceitual, em linha com margens praticadas nessa fase~\cite{r3}.

\begin{table}[ht]
\centering
\caption{Orçamento de massa e potência (cenários paramétricos, não validados).}
\label{tab:massa}
\small
\begin{tabular},{@{}lcc@{}}
\toprule
Item & Caso A (LTAN 06:00) & Caso B (LTAN $\approx$ 12:00)\\
\midrule
Eclipse máximo (min)                            & 18,8 & 35,3\\
Energia de eclipse (kWh)                        & 328,6 & 617,5\\
Bateria (kWh / kg)                              & 939,0 / 5.691 & 1.764,3 / 10.693\\
Potência de geração (MW)                        & 1,359 & 1,736\\
Área do arranjo (m$^2$) / massa (kg)            & 4.373 / 13.992 & 5.588 / 17.882\\
Área radiadora em planta (m$^2$) / massa (kg)   & 1.528 / 9.170 & 1.550 / 9.301\\
TI (kg)                                         & 6.000 & 6.000\\
Estrutura e chicotes (kg)                       & 5.228 & 6.581\\
ADCS (kg)                                       & 2.746 & 5.496\\
Hardware diverso (kg)                           & 300 & 300\\
\textbf{Massa seca com margem (kg)}             & \textbf{51.752} & \textbf{67.504}\\
Propelente de arrasto (kg)                      & 636 & 812\\
\textbf{Massa de lançamento (kg)}               & \textbf{52.388} & \textbf{68.316}\\
\bottomrule
\end{tabular}
\end{table}

\paragraph{Lançamentos.} Sob a hipótese de 20--25~t por voo ao SSO com
recuperação parcial (hipótese sem fonte primária), seriam 3 lançamentos (A) e
3--4 (B). As capacidades publicadas pelo fabricante (22.800~kg e 63.800~kg para
LEO nos Falcon~9 e Falcon~Heavy; mais de 100~t projetadas para a
Starship~\cite{r10,r11,r12}) valem para outras órbitas e configurações e não
garantem entrega em SSO. Montagem em órbita e volume de carenagem das mantas não
foram avaliados.

\subsection{Confiabilidade e redundância}\label{sec:conf}

Toma-se como taxa-base a observada no treinamento do Llama~3: 419 interrupções
inesperadas em 54 dias com até 16.384 GPUs H100~\cite{r36}. Isso resulta em
$\lambda=1{,}97\times10^{-5}$~h$^{-1}$ por GPU, tratando todas as interrupções
como falhas (limite superior: cerca de 78\% foram atribuídas a
hardware~\cite{r36}). Seja $\varphi$ a fração de interrupções que são falhas
permanentes (P15: 5--30\%). Com modelo exponencial independente, sem reparo:
\begin{equation}
  p_s=e^{-\varphi\lambda t},\qquad
  n_{min}=\min\Big\{n:\ \Pr[\mathrm{Bin}(n,p_s)\ge N_{req}]\ge0{,}99\Big\},\quad
  N_{req}=1000,\ t=87.600~\text{h}.
  \label{eq:rel}
\end{equation}

\begin{table}[ht]
\centering
\caption{Sobredimensionamento necessário (modelo exploratório).}
\label{tab:rel}
\small
\begin{tabular},{@{}cccc@{}}
\toprule
$\varphi$ & $p_s$ (10 anos) & $n_{min}$ & Sobredimensionamento\\
\midrule
5\%  & 0,917 & 1.114 & 11\%\\
10\% & 0,841 & 1.225 & 23\%\\
20\% & 0,708 & 1.471 & 47\%\\
30\% & 0,595 & 1.760 & 76\%\\
\bottomrule
\end{tabular}
\end{table}

\emph{Limitação:} a taxa terrestre não inclui radiação (SEE, TID) e deve ser um
limite inferior para o ambiente orbital~\cite{r48}. Para uma vida de GPUs COTS
em LEO de cerca de 3 anos~\cite{r19}, uma missão de dez anos exigiria até três
sobressalentes por unidade ativa; o sobredimensionamento da
Tabela~\ref{tab:rel} deve ser lido como limite inferior.

\subsection{Barramento, tempo discreto e requisitos de GNC e estimação}

Com barramento de 400~V (hipótese sem fonte primária) e 1,05~MW, a corrente é
$\approx 2{,}6$~kA. Cargas de IA comportam-se como cargas de potência constante,
com impedância incremental negativa $Z_{in}=-V^{2}/P=-0{,}152~\Omega$. A
estabilidade do barramento CC exige que a impedância de saída da fonte
satisfaça $|Z_{out}|\ll|Z_{in}|$ em todas as frequências (critério de
Middlebrook~\cite{r34}; ver também~\cite{r33}). Arranjos de alta tensão em
plasma de LEO exigem análise de arco e de interação com plasma~\cite{r35}.

Regras de projeto, apresentadas como requisitos e não como resultados:
\begin{itemize}
  \item A taxa de amostragem de malhas de controle discretas deve ser 10 a 20
  vezes a largura de banda em malha fechada~\cite{r32}.
  \item A largura de banda de atitude deve ficar bem abaixo do primeiro modo
  flexível dos arranjos ou usar filtros \emph{notch}/estabilização de
  fase~\cite{r27}.
  \item A estimação de atitude exige filtragem do tipo EKF multiplicativo e
  sensores com modelo de ruído~\cite{r31,r49}.
  \item A estimação de SOC e SOH da bateria e da temperatura dos radiadores é
  necessária, pois esses estados não são medidos diretamente.
\end{itemize}
Nada disso foi simulado neste trabalho.

% ---------------------------------------------------------------------
\section{Inventário de ciclo de vida}

\subsection{Hardware avançado}

A pegada de aceleradores depende de nó tecnológico, rendimento, área, memória
HBM e encapsulamento. Não existe fator universal por \emph{chip}, e a pegada de
um processador de 3~nm não pode ser inferida do nó nominal. O inventário deve
vir de fabricantes, declarações ambientais de produto ou bases com
representatividade documentada.
\begin{equation}
  G_{electr}=\sum_{i=1}^{n} m_i\,g_i,
  \label{eq:gelectr}
\end{equation}
em que $m_i$ e $g_i$ devem estar em unidades compatíveis (massa $\times$
impacto por massa; unidades $\times$ impacto por unidade; atividade $\times$
impacto por atividade).

A cadeia material depende de minerais críticos (gálio, germânio, índio, lítio,
cobre)~\cite{r21}. Eles entram como categoria de impacto a ser selecionada e
justificada nos termos da Seção~\ref{sec:norma}.

\subsection{Lançamentos}\label{sec:lanc}

Para o conjunto de $N_{launch}$ voos, com lançadores parcialmente reutilizáveis:
\begin{equation}
  G_{launch}=\sum_{j=1}^{N_{launch}}\left[
    \frac{G_{mfg,boost,j}},{N_{reuse,j}}+G_{mfg,exp,j}
    +m_{prop,j}\,I_{prop,j}+G_{ground,j}\right],\qquad
  g_{launch}=\frac{\sum_j \alpha_j\,G_{launch,j}},{m_{payload,delivered}}.
  \label{eq:launch}
\end{equation}
$G_{mfg,exp,j}$ é a fabricação dos elementos descartáveis (estágio superior,
carenagem); $I_{prop,j}$ é o fator de ciclo de vida do propelente
(kg~$\COeq$/kg); $\alpha_j$ é a fração alocada à missão (massa, volume, energia
ou causalidade, com sensibilidade). O denominador é a massa entregue à órbita
de destino.

\paragraph{Piso físico de $\CO$ de combustão.} Pela equação de
Tsiolkovsky~\cite{r30}, com massa estrutural nula (limite inferior de propelente
por kg de carga) e razão de mistura fixa:
\begin{equation}
  \frac{m_{\CO}},{m_{payload}}\;\ge\;
  \left(e^{\Delta v/(g_0 I_{sp})}-1\right)x_f\,c_{\CO}.
  \label{eq:floor}
\end{equation}
Aqui $x_f$ é a fração de combustível (1/3,6 para RP-1/LOX a 2,6:1; 1/4,6 para
CH$_4$/LOX a 3,6:1) e $c_{\CO}$ é 3,15~kg/kg (RP-1, 86\% de carbono) ou
2,74~kg/kg (CH$_4$). Com $\Delta v=9{,}0$--$9{,}5$~km/s e $I_{sp}$ efetivo de
330--365~s (RP-1) ou 350--380~s (CH$_4$) (P17), o piso é de
\textbf{9,9--15,6~kg $\CO$/kg} (RP-1) e \textbf{6,1--8,9~kg $\CO$/kg}
(CH$_4$). É um piso, não uma estimativa: ignora estrutura, recuperação,
produção de propelente e operações de solo. Para comparação, os termos de
propelente de~\cite{r18} resultam em 28,0~kg $\CO$/kg de carga (Falcon~9:
155,87~t de RP-1 para 17,5~t de carga) e 18,9~kg $\CO$/kg (Starship: 1.030~t de
CH$_4$ para 150~t de carga), acima do piso, como esperado, pois refletem a massa
estrutural e as perdas reais dos veículos.

\subsection{Emissões atmosféricas e reentrada}\label{sec:atm}

Maloney \emph{et al.} modelaram um cenário-base de 10~Gg/ano de carbono negro e
cenários de 30 e 100~Gg/ano~\cite{r13}. Ryan \emph{et al.} desenvolveram um
inventário de lançamentos e reentradas com efeitos modelados em ozônio e
clima~\cite{r14}. Brown \emph{et al.} publicaram o inventário global de 2019,
que estima, para a estratosfera (15--50~km), 5,82~Gg de $\CO$, 6,38~Gg de vapor
d'água, 0,28~Gg de carbono negro, 0,22~Gg de NO$_x$, 0,50~Gg de cloro reativo e
0,91~Gg de alumina~\cite{r15}. Esses valores são inventários globais, não
fatores por veículo.

Ferreira \emph{et al.} estimam que um satélite típico de 250~kg pode gerar cerca
de 30~kg de nanopartículas de alumina na reentrada, e que satélites reentrando
em 2022 produziram cerca de 17~t de compostos de alumínio. Em cenários de
megaconstelações, a produção poderia superar cerca de 360~t/ano, com
implicações para o ozônio~\cite{r16}. Outras espécies (cobre, lítio) e efeitos
sobre o albedo estratosférico carecem de quantificação nas fontes consultadas. A
composição de um ODC (compósitos de carbono, silício, GaAs, baterias de íons de
lítio) difere de um satélite típico, e esse resultado não é extrapolável
diretamente.

Ohs \emph{et al.}~\cite{r18} apresentam a ferramenta ESpaS para estimar
intensidades de carbono de CPU, memória e rede em cenários orbitais e
terrestres, com três exemplos de aplicação. Para sistemas dimensionados para
cinco anos de missão, estimam intensidades de lançamento e reentrada de
51,7~kg~$\COeq$/kg de carga (Falcon~9, carga de 17,5~t) e de 34,4~kg~$\COeq$/kg
(Starship, melhor caso de 150~t). A reentrada contribui com 13,2~kg~$\COeq$/kg
de massa reentrante, a partir de um fator de 0,4~kg NO$_x$/kg convertido em
$\COeq$, e a alumina não é contabilizada numericamente por falta de conversão
para $\COeq$. Incluindo a bateria necessária ao eclipse, a intensidade de
energia resulta em 34,0 (Terra), 66,3 (Falcon~9) e 52,1~g~$\COeq$/kWh
(Starship), e a intensidade de CPU em carga plena em 282,8, 566,3 e
445,4~$\mu$g~$\COeq$/s, ou seja, de 1,5 a 2,0 vezes a terrestre. Os autores
concluem que, mesmo sob hipóteses otimistas, sistemas em órbita têm custo de
carbono significativamente maior, principalmente por emissões incorporadas de
lançamento e reentrada. A versão v3 (fev/2026) corrige constantes de efeitos
químicos, e o trabalho foi publicado no \emph{ACM SIGENERGY Energy Informatics
Review}~\cite{r18}. Os valores vêm de análises do tipo
\emph{back-of-the-envelope}, nas palavras dos autores, aplicadas a uma
configuração de exemplo de pequeno porte (arranjo de 2~m$^2$, bateria de
4~kWh, CPU de 28~W). Por isso, constituem evidência favorável à cautela, não
prova numérica para o ODC aqui modelado.

\paragraph{Comparação com o ESpaS-ODC~\cite{r19}.} Ohs \emph{et al.}
propõem o ESpaS-ODC, um modelo aberto de carbono de ciclo de vida para
\emph{data centers} orbitais, aplicado a um nó de 1~kW (uma GPU H100) a 510~km.
O modelo inclui três subsistemas que os autores consideram fisicamente
inevitáveis em densidades de potência de classe GPU: um radiador térmico
calibrado com dados da ISS, a degradação do arranjo solar com margem de eclipse
e sobressalentes em \emph{standby} frio, dada a ausência de acesso para reparo.
A métrica é o carbono por GPU-hora,
$I_{\text{GPU-h}}=(I_{L\&R}\,m_{total}+C_{prod})/T$, com $I_{L\&R}$
de~\cite{r18}. O barramento (estrutura, ADCS, aviônica) fica fora do modelo, e a
fronteira terrestre considera apenas o carbono incorporado do \emph{chip}, o
que os autores descrevem como comparação conservadora.

Nesse varredor, a paridade com um \emph{data center} terrestre de média global
(477~g~$\COeq$/h; PUE 1,54) ocorre em $T\le1{,}0$~ano (Starship) e
$T\approx2{,}0$~anos (Falcon~9). Frente a um \emph{data center} ``verde''
(107~g~$\COeq$/h; rede finlandesa de 68~g/kWh; PUE 1,2), a paridade exige
$T=6$~anos e duas unidades (Starship) ou $T=9$~anos e três (Falcon~9). Num nó
de 1~kW, o radiador responde por 30,5~kg de 44,3~kg de componentes modelados, e
um sobressalente frio aumenta o carbono por GPU-hora em 40\% (Starship) e 34\%
(Falcon~9) a $T=3$~anos~\cite{r19}. Os próprios autores afirmam não poder dizer
qual opção é melhor, mas contestam a tese de que \emph{data centers} orbitais
seriam a única solução sustentável~\cite{r19}.

Este trabalho difere de~\cite{r19} em três pontos: modela um nó de 1~MW com
barramento, ADCS e eclipse por LTAN; trata água e tetos de emissão incorporada;
e propaga incerteza. Em contrapartida, não calcula carbono incorporado nem
carbono por GPU-hora. Dois parâmetros de~\cite{r19} merecem comparação direta.
Primeiro, o radiador de~\cite{r19} resulta em
$3{,}03~\text{m}^2/\text{kW}\times10~\text{kg/m}^2\approx30{,}3$~kg/kW, contra
$\approx8$--$9$~kg/kW neste trabalho (6~kg/m$^2$ de área em planta com fluxo
bifacial de 360~W/m$^2$). Na parametrização aqui adotada, isso equivale a
$\rho_{rad}\approx21{,}8$~kg/m$^2$, valor fora da faixa de 4--9~kg/m$^2$ usada
inicialmente e que motivou a ampliação para 4--22~kg/m$^2$. A fração do
radiador (17,7\% da massa seca no caso A) não é comparável diretamente com a de
68,8\% de~\cite{r19}, porque os denominadores diferem: massa de componentes de
um nó de 1~kW sem barramento, de um lado, e massa seca do nó de 1~MW com
estrutura, ADCS e margem, de outro. Segundo, \cite{r19} indica que GPUs COTS
falham em cerca de 3~anos em LEO, o que implica $r\ge\lceil T/3\rceil$; para
$T=10$~anos, $r\ge4$, isto é, até três sobressalentes por unidade ativa, ante
11\%--76\% de sobredimensionamento na Tabela~\ref{tab:rel}. Essa extrapolação
da regra de~\cite{r19} para dez anos é dos autores deste trabalho, não
de~\cite{r19}.

\begin{table}[ht]
\centering
\caption{Massa de lançamento (t) ao substituir parâmetros do modelo pelos
de~\cite{r19}. Sensibilidade, não validação.}
\label{tab:cal19}
\small
\begin{tabular},{@{}lcc@{}}
\toprule
Cenário & Caso A & Caso B\\
\midrule
Caso-base                                  & 52,4 & 68,3\\
Manta: 3,41~kg/m$^2$                       & 53,9 & 70,4\\
Bateria: 3,75~kg/kWh                       & 49,4 & 62,7\\
Radiador: 30,3~kg/kW                       & 85,7 & 102,2\\
Manta, bateria e radiador de~\cite{r19}    & 84,2 & 98,6\\
\bottomrule
\end{tabular}
\end{table}

Com o radiador de~\cite{r19}, a massa de lançamento fica acima das medianas do
Monte Carlo (71,0~t no caso A e 89,3~t no caso B) e dentro dos intervalos de
90\% (57,2--91,3~t e 73,0--111,1~t), o que decorre de a faixa de $\rho_{rad}$
ter sido ampliada justamente para abrigar esse valor. Aplicando as
intensidades de lançamento e reentrada de~\cite{r18} (34,4--51,7~kg~$\COeq$/kg)
às massas da Tabela~\ref{tab:cal19} com os três parâmetros de~\cite{r19}, esses
termos ocupam de 4,9\% a 10,8\% do teto em rede de 450~g/kWh e de 44\% a 97\%
em rede de 50~g/kWh. Portanto, a conclusão qualitativa da
Seção~\ref{sec:teto} se mantém, mas a margem em redes de baixa intensidade é
sensível à massa do radiador.

\subsection{Linha de base terrestre}
\begin{equation}
  E_{facility}=E_{IT}\cdot PUE,\qquad
  G_{elec}=E_{facility}\,I_{grid},\qquad
  W_{direct}=E_{IT}\cdot WUE.
  \label{eq:base}
\end{equation}
PUE e WUE seguem a ISO/IEC 30134-2 e 30134-9~\cite{r41,r42}. A intensidade
$I_{grid}$ deve ser regional e temporal, ou marginal quando a decisão depende
de carga adicional. Como referência de ordem de grandeza, o
Öko-Institut~\cite{r8} adota 0,393 (EUA), 0,584 (China), 0,300 (Europa) e
0,484~kg~$\COeq$/kWh (demais regiões), a partir de dados do Ember. Para WUE,
\cite{r8}, com dados do LBNL (2024), reporta 0,32--0,67~L/kWh por categoria de
\emph{data center} dos EUA (média de 0,36) e relatos de operadores de até
5~L/kWh. Revisões sobre consumo de \emph{data centers} estão em~\cite{r40}.

% ---------------------------------------------------------------------
\section{Resultados e discussão}

\subsection{Energia e água de referência}
\begin{equation}
  t=87.600~\text{h},\quad E_{IT}=87{,}6~\text{GWh},\quad
  E_{facility}=105{,}12\text{--}131{,}4~\text{GWh (PUE 1,2--1,5)}.
  \label{eq:energia}
\end{equation}
São resultados aritméticos da UF, sem indisponibilidade, degradação ou
substituição.

\begin{table}[ht]
\centering
\caption{Consumo direto de água da linha de base terrestre em dez anos
(milhões de litros).}
\label{tab:agua}
\small
\begin{tabular},{@{}lcc@{}}
\toprule
Cenário de WUE & Faixa (L/kWh de TI) & $W_{direct}$\\
\midrule
Caso-base (LBNL/Öko~\cite{r8})     & 0,32--0,67 & 28,0--58,7\\
Cenário dos autores (ver texto)    & 1--2       & 87,6--175,2\\
Extremo reportado por operadores   & 5          & 438\\
\bottomrule
\end{tabular}
\end{table}

\paragraph{Água indireta da geração.} Aplicando os fatores de~\cite{r8} à
eletricidade da instalação (105,12--131,4~GWh): 0,28~L/kWh (eólica) resulta em
29,4--36,8 milhões de litros; 2,33~L/kWh (carvão), em 244,9--306,2 milhões;
$\sim$70~L/kWh (hidrelétrica), em 7.358--9.198 milhões. Em relação ao caso-base
direto (28,0--58,7 milhões), isso equivale a cerca de 0,5 a 1,3 vez (eólica), 4
a 11 vezes (carvão) e 125 a 330 vezes (hidrelétrica). A economia direta de um
ODC, se houver, é assim uma fração pequena ou grande da pegada hídrica total
conforme a matriz elétrica. Além disso, litros de regiões de escassez não
equivalem a litros de regiões de abundância~\cite{r39}, e a pegada hídrica da
fabricação de \emph{chips}, painéis e baterias, e da produção de propelentes,
deve ser incluída.

\paragraph{Verificação de ordem de grandeza do WUE.} O Jornal da
Unesp~\cite{r22}, atribuindo o dado à IEA, informa consumo global de água por
\emph{data centers} de cerca de 560 bilhões de litros por ano, com projeção de
1,2 trilhão até 2030. Dividido pelos $\approx$415~TWh de~\cite{r5}, isso
equivale a $\approx$1,35~L/kWh de eletricidade das instalações, ou
$\approx$1,6--2,0~L/kWh de energia de TI para PUE de 1,2--1,5. Esse valor está
na faixa do cenário dos autores (1--2~L/kWh) e acima da faixa de
0,32--0,67~L/kWh das categorias dos EUA~\cite{r8}. É uma estimativa de ordem de
grandeza: depende de um relato jornalístico, de possível diferença de ano-base
entre os dois números e da conversão por PUE, e não é uma medição de WUE.

\subsection{Balanço de carbono e \emph{payback}}
\begin{align}
  G_{ODC}&=G_{mfg}+G_{launch}+G_{op}+G_{repl}+G_{EOL},\nonumber\\
  G_{ground}&=G_{mfg,g}+G_{elec}+G_{cool}+G_{repl,g}+G_{EOL,g},
  \label{eq:bal}\\
  \Delta G&=G_{ground}-G_{ODC}.
  \label{eq:dg}
\end{align}
Se $\Delta G>0$, o ODC tem menor GWP dentro das fronteiras e hipóteses
adotadas. Define-se $\Delta G_{inc}=G_{inc,ODC}-G_{inc,terr}$.

Com redução operacional anual constante e $G_{ODC,yr}$ as emissões
operacionais anuais do ODC:
\begin{equation}
  t_{payback}=\frac{\Delta G_{inc}},{E_{IT,yr}\,PUE\,I_0-G_{ODC,yr}},\qquad
  E_{IT,yr}=8{,}76\times10^{6}~\text{kWh/ano}.
  \label{eq:payback}
\end{equation}
Para $I_{grid}(t)=I_0e^{-\gamma t}$ e $A=E_{IT,yr}\,PUE\,I_0$, o \emph{payback}
satisfaz
\begin{equation}
  \Delta G_{inc}+G_{ODC,yr}\,t_p-\frac{A},{\gamma}\left(1-e^{-\gamma t_p}\right)=0,
  \label{eq:paybackexp}
\end{equation}
que é transcendental para $G_{ODC,yr}\neq0$. Para $G_{ODC,yr}=0$:
\begin{equation}
  t_p=-\frac{1},{\gamma}\ln\!\left(1-\frac{\gamma\,\Delta G_{inc}},{A}\right),
  \qquad \frac{\gamma\,\Delta G_{inc}},{A}<1.
  \label{eq:paybackclosed}
\end{equation}
Se $A/\gamma<\Delta G_{inc}$, não há \emph{payback} em horizonte infinito. As
reposições ($G_{repl}$) ocorrem ao longo do tempo. Ao agrupá-las em
$\Delta G_{inc}$ (\emph{upfront}), o \emph{payback} fica enviesado a favor do
terrestre. Como as fontes públicas não permitem um inventário comparável,
\textbf{não se apresenta \emph{payback} numérico}.

\subsection{Teto de emissões incorporadas diferenciais}\label{sec:teto}

Como $\Delta G_{inc}$ é desconhecido, a Tabela~\ref{tab:teto} fornece o teto
$G_{inc}^{max}=E_{facility,10y}\,I_{grid}$, o valor máximo que a diferença de
emissões incorporadas pode atingir para que a economia operacional terrestre
compense o investimento ambiental em dez anos. Ele carrega dois vieses de
sinais opostos. Ao fixar $G_{ODC,yr}=0$ (omite estações em terra, enlaces e
contingências), o teto fica mais favorável ao ODC. Ao omitir as emissões
incorporadas terrestres, fica mais desfavorável. Os dois efeitos não estão
quantificados e não se compensam por construção.

\begin{table}[ht]
\centering
\caption{Teto de $G_{inc}$ para PUE 1,2--1,5, em dez anos. Intensidades
paramétricas, não regionais.}
\label{tab:teto}
\small
\begin{tabular},{@{}ccccc@{}}
\toprule
$I_{grid}$ & Anual & 10 anos & kg $\COeq$/kg, A & kg $\COeq$/kg, B\\
(g $\CO$/kWh) & (t $\COeq$) & (t $\COeq$) & (52,4~t) & (68,3~t)\\
\midrule
450 & 4.730--5.913 & 47.304--59.130 & 903--1.129 & 692--866\\
350 & 3.679--4.599 & 36.792--45.990 & 702--878   & 539--673\\
220 & 2.313--2.891 & 23.126--28.908 & 441--552   & 339--423\\
50  & 526--657     & 5.256--6.570   & 100--125   & 77--96\\
\bottomrule
\end{tabular}
\end{table}

O valor por kg é uma métrica de projeto (intensidade incorporada por massa
lançada) e não um fator de emissão do lançador. Em rede de baixa intensidade
(50~g/kWh), o teto cai para 77--125~kg $\COeq$ por kg lançado.

\paragraph{Piso de combustão versus teto.} O piso da Seção~\ref{sec:lanc}
(6,1--15,6~kg $\CO$/kg) corresponde a $\le 15{,}5\%$ (A) e $\le 20{,}3\%$ (B) do
menor teto da Tabela~\ref{tab:teto} (5.256~t). Logo, a combustão, como piso, não
esgota o teto.

\paragraph{Verificação cruzada com~\cite{r18}.} Multiplicando as intensidades
de lançamento e reentrada de~\cite{r18} (51,7~kg $\COeq$/kg para Falcon~9; 34,4
para Starship) pelas massas de lançamento do caso-base (52,4~t e 68,3~t),
obtêm-se 2.708--3.532~t $\COeq$ (Falcon~9) e 1.802--2.350~t $\COeq$ (Starship).
Em relação ao teto, isso é de 3,0\% a 7,5\% em rede de 450~g/kWh e de 27,4\% a
67,2\% em rede de 50~g/kWh. Com as massas do P95 do Monte Carlo
(Seção~\ref{sec:mc}), os valores sobem para até 10,0\% (A) e 12,1\% (B) em rede
de 450~g/kWh, e para até 89,8\% (A) e 109,3\% (B) em rede de 50~g/kWh; no P95 do
caso B, com Falcon~9, PUE 1,2 e 50~g/kWh, esses dois termos excederiam o menor
teto em cerca de 9\%. Ressalvas: as intensidades de~\cite{r18} dependem de
cargas de 17,5~t (Falcon~9) e 150~t (melhor caso do Starship), do fator de
0,4~kg NO$_x$/kg na reentrada e da extrapolação linear da fabricação do
Starship, adotados naquele trabalho; excluem a alumina e a fabricação do ODC.
Assim, em redes de baixa intensidade lançamento e reentrada podem ocupar parcela
relevante do teto, enquanto em redes de alta intensidade são secundários frente
à economia operacional terrestre. Não se infere vantagem do ODC a partir disso.

\subsection{Sensibilidade e incerteza}\label{sec:mc}

A variação um-por-vez (OAT) sobre a massa de lançamento é resumida na
Tabela~\ref{tab:oat} e na Figura~\ref{fig:tornado}, que traz os oito parâmetros
que aparecem entre os sete de maior efeito em pelo menos um dos casos.
Parâmetros de menor efeito (por exemplo, $q$ do radiador e $\rho_{atm}$) foram
omitidos da tabela.

\begin{table}[ht]
\centering
\caption{Sensibilidade OAT da massa de lançamento (t), ordenada pelo efeito no
caso A.}
\label{tab:oat}
\small
\begin{tabular},{@{}lccc@{}}
\toprule
Parâmetro & Faixa testada & Faixa em A (t) & Faixa em B (t)\\
\midrule
$\rho_{rad}$ (kg/m$^2$) & 4,0--22,0  & 48,2--86,1 & 64,0--102,5\\
$\rho_{PV}$ (kg/m$^2$)  & 2,5--4,5   & 47,4--61,6 & 61,5--81,0\\
$M_{IT}$ (kg/kW)        & 6--12      & 52,4--60,7 & 68,3--76,6\\
$k_{arr}$               & 0,70--0,90 & 49,2--56,7 & 63,6--75,0\\
$k_{st}$                & 0,10--0,25 & 50,3--56,6 & 65,7--73,6\\
DoD                     & 0,25--0,50 & 50,0--55,5 & 63,9--74,2\\
$\eta_{cell}$           & 0,27--0,32 & 50,2--55,2 & 65,0--72,6\\
$e_b$ (Wh/kg)           & 120--200   & 51,0--55,3 & 65,7--73,9\\
\bottomrule
\end{tabular}
\end{table}

No caso A, os cinco maiores efeitos são $\rho_{rad}$, $\rho_{PV}$, $M_{IT}$,
$k_{arr}$ e $k_{st}$. No caso B, são $\rho_{rad}$, $\rho_{PV}$, $k_{arr}$, DoD e
$M_{IT}$. O DoD e a densidade de bateria pesam mais no caso B, devido ao
eclipse longo.

\begin{figure}[ht]
\centering
\figinc{fig_tornado.pdf},{\textwidth}
\caption{Diagrama de tornado da sensibilidade OAT da massa de lançamento nos
casos A e B. Cada barra parte da massa-base e mostra o efeito do valor mínimo
(azul) e do valor máximo (vermelho) da faixa do parâmetro.}
\label{fig:tornado}
\end{figure}

O método de Monte Carlo com distribuições triangulares independentes (P18;
$N=3000$, semente 1; modos iguais ao caso-base, exceto $M_{IT}=8$~kg/m$^2$ e
$N_{red}=2$) produz P5/P50/P95 de \textbf{57,2 / 71,0 / 91,3~t} (A) e
\textbf{73,0 / 89,3 / 111,1~t} (B). Os P5 e P50 acima do caso-base decorrem da
assimetria das faixas (modos deslocados em relação ao ponto médio). As faixas
são hipóteses, não dados.

O resultado depende da distribuição adotada para $\rho_{rad}$, o parâmetro de
maior efeito. A Tabela~\ref{tab:mcvar} mostra cinco variantes. A mediana varia
de 62,1 a 85,0~t no caso A e de 79,8 a 102,8~t no caso B, e a escolha da
distribuição é subjetiva. Por isso, os percentis do Monte Carlo ilustram a
propagação de incerteza e não devem ser lidos como probabilidades; devem ser
refeitos com distribuições sustentadas por dados.

\begin{table}[ht]
\centering
\caption{Sensibilidade do Monte Carlo à distribuição de $\rho_{rad}$ (massa de
lançamento, t; P5 / P50 / P95).}
\label{tab:mcvar}
\small
\begin{tabular},{@{}lcc@{}}
\toprule
Distribuição de $\rho_{rad}$ & Caso A & Caso B\\
\midrule
Triangular 4--9, moda 6           & 53,7 / 62,1 / 72,0 & 68,7 / 79,8 / 93,3\\
Triangular 4--22, moda 6 (base)   & 57,2 / 71,0 / 91,3 & 73,0 / 89,3 / 111,1\\
Triangular 4--22, moda 12         & 61,2 / 76,4 / 93,8 & 77,4 / 94,4 / 113,9\\
Triangular 4--22, moda 21,8       & 64,8 / 85,0 / 101,1 & 81,4 / 102,8 / 121,5\\
Uniforme 4--22                    & 57,5 / 77,4 / 98,7 & 73,8 / 95,4 / 118,8\\
\bottomrule
\end{tabular}
\end{table}

\subsection{Pegada por capacidade computacional}

A comparação por MW instalado pode ocultar diferenças de desempenho útil
(redundância, limite térmico, comunicação, disponibilidade). A
Eq.~\eqref{eq:igwp} exige métrica $Q$ declarada: modelo de IA de referência,
precisão numérica, lote, utilização, latência, disponibilidade, memória e
transferência de dados. Pico nominal de TFLOPS não representa fluxo de trabalho
equivalente.

\subsection{Principais gargalos candidatos}

Fabricação de aceleradores, HBM e encapsulamento; painéis e estruturas
térmicas; baterias e conversores; propelentes e lançadores; redundância e
substituição; reentrada. A ordenação depende de inventários reais. Nem o
silício nem o lançamento podem ser declarados dominantes sem dados.

% ---------------------------------------------------------------------
\section{Recomendações}

\subsection{Transparência de inventário}
Operadoras devem divulgar:
\begin{enumerate}
  \item massa seca lançada e massa entregue;
  \item órbita, lançador, configuração e perfil da missão;
  \item propelentes, emissões e regra de alocação;
  \item massa e origem de semicondutores, painéis e baterias;
  \item energia incorporada e dados de fabricação;
  \item capacidade computacional útil e disponibilidade;
  \item degradação, substituições e redundância;
  \item consumo direto e indireto de água;
  \item destino final e massa estimada de produtos de reentrada;
  \item avaliações de impacto ambiental cumulativo exigidas por agências de
  licenciamento;
  \item incertezas, limitações e resultados de revisão crítica.
\end{enumerate}

\subsection{Requisitos de engenharia para um estudo de segunda fase}
\begin{itemize}
  \item \textbf{GNC:} simulação 6-DOF com modos flexíveis e gestão de momento
  (TEA)~\cite{r27,r31}, com propagação orbital conforme~\cite{r28}.
  \item \textbf{Atmosfera e detritos:} modelo atmosférico~\cite{r29} e
  conformidade com requisitos de mitigação de detritos~\cite{r43}.
  \item \textbf{Ambiente espacial:} especificação conforme~\cite{r44}
  e~\cite{r48}.
  \item \textbf{Térmica, AOCS e fotovoltaico:} normas
  ECSS-E-ST-31C~\cite{r45}, ECSS-E-ST-60-30C~\cite{r46} e
  ECSS-E-ST-20-08C~\cite{r47}.
  \item \textbf{Qualificação:} TID/SEE, ciclagem térmica de $10^{4}$--$10^{5}$
  ciclos e qualificação de baterias em ciclagem rasa de longa duração.
  \item \textbf{Sistemas:} ConOps, orçamento de enlace e de \emph{jitter},
  montagem em órbita e análise de falhas.
\end{itemize}

\subsection{Mitigação}
Reduzir massa por unidade de computação útil, modularidade, extensão de vida,
redução de substituições, reparabilidade e política de fim de vida com
reentrada controlada ou órbita cemitério avaliada por impacto atmosférico,
risco de colisão, energia e propelente.

% ---------------------------------------------------------------------
\section{Conclusões}

\emph{Data centers} orbitais não podem ser considerados ambientalmente
superiores apenas por usarem energia solar extraterrestre ou dispensarem torres
evaporativas. Para 1~MW de TI por dez anos, a energia de TI é 87,6~GWh e a
eletricidade total de uma linha de base terrestre é 105,12--131,4~GWh (PUE
1,2--1,5). O consumo direto de água é de 28,0--58,7 milhões de litros no
caso-base de WUE, mas a água da geração elétrica é de ordem comparável (matriz
eólica) a mais de duas ordens de grandeza maior (matriz hidrelétrica).

O modelo paramétrico do nó orbital resulta em 52,4~t (A) e 68,3~t (B) de massa
de lançamento, com intervalos de 90\% de 57,2--91,3~t e 73,0--111,1~t. A massa
areal do radiador, a das mantas e a massa específica da eletrônica dominam a
incerteza, e o intervalo depende da distribuição adotada para a massa areal do
radiador. Sob a premissa adotada, falhas de hardware exigem 11\% a 76\% de
sobredimensionamento. Aplicando as intensidades de lançamento e reentrada
de~\cite{r18}, esses termos ocupam de 3\% a 7,5\% do teto de emissões
incorporadas em rede de 450~g/kWh e de 27\% a 67\% em rede de 50~g/kWh, e
podem excedê-lo no P95 do caso B.

A vantagem ambiental líquida permanece condicional. Ela só poderá ser
demonstrada com inventário específico, unidade funcional de computação útil,
alocação transparente de lançamentos, avaliação de incerteza e revisão crítica.

\section*{Limitações e disponibilidade de dados}
Todos os valores derivam de parâmetros declarados (Tabela~\ref{tab:premissas})
e não de medições. O modelo orbital é de primeira ordem. Faltam modos flexíveis,
radiação, MMOD, penumbra e simulação 6-DOF. As taxas de falha vêm de um
\emph{cluster} terrestre. O Monte Carlo usa faixas assumidas e depende da
distribuição de $\rho_{rad}$. Os tetos da Tabela~\ref{tab:teto} omitem as
emissões operacionais do ODC e as incorporadas terrestres.

O código de cálculo (\texttt{odc\_lca\_model.py}), que inclui verificação
automática contra os valores deste artigo, e o de geração das figuras
(\texttt{odc\_lca\_figures.py}) estão disponíveis em
\url{https://github.com/dafnhe-thor-barros/Article07}.
% Antes de submeter: arquivar uma versao no Zenodo e citar o DOI aqui.

\section*{Declaração de conflito de interesses}
Os autores declaram ausência de relações comerciais, financeiras, contratuais ou
institucionais que possam ser interpretadas como conflito de interesses, sem
vínculo com fabricantes de lançadores, fornecedores de hardware ou operadoras
de \emph{data centers}.

% ---------------------------------------------------------------------
\appendix
\section{Premissas e origem}

\begin{longtable},{@{}p{0.07\textwidth}p{0.25\textwidth}p{0.26\textwidth}p{0.30\textwidth}@{}}
\caption{Premissas do modelo e sua origem.}\label{tab:premissas}\\
\toprule
ID & Parâmetro & Valor & Origem\\
\midrule
\endfirsthead
\toprule
ID & Parâmetro & Valor & Origem\\
\midrule
\endhead
\bottomrule
\endfoot
P1  & $P_{load}$ & 1,05~MW & Hipótese (TI + 5\%); sem fonte primária\\
P2  & $\eta_{conv}$, $\eta_{rt}$ & 0,97; 0,92 & Hipótese; sem fonte primária\\
P3  & DoD, $e_b$ & 0,35; 165~Wh/kg & Hipótese; sem fonte primária\\
P4  & $\eta_{cell}$ (BOL), $k_{arr}$ & 0,295; 0,80 & Hipótese; sem fonte primária\\
P5  & $\rho_{PV}$ & 3,2~kg/m$^2$ (2,5--4,5) & Hipótese; validar com \emph{datasheet}\\
P6  & $q$, $\rho_{rad}$ & 360~W/m$^2$; 6,0~kg/m$^2$ (4--22) & Hipótese; faixa de $\rho_{rad}$ cobre o equivalente de~\cite{r19} (21,8)\\
P7  & $M_{IT}$ & 6~kg/kW (6--12) & Hipótese; sem fonte primária\\
P8  & Estrutura e chicotes & 15\% (10--25\%) & Hipótese; faixa de projeto conceitual~\cite{r26}\\
P9  & Margem de massa & 20\% & Hipótese; prática de fase conceitual~\cite{r3}\\
P10 & $h_{spec}$, $N_{red}$ & 17,5~N\,m\,s/kg; 2 & Hipótese; validar com \emph{datasheet}\\
P11 & Geometria dos arranjos & 2 asas de 30~m; cubo de 10~m & Hipótese\\
P12 & $M_{hw}$ & 300~kg & Hipótese\\
P13 & Arrasto & $\rho=10^{-13}$~kg/m$^3$; $C_d=2{,}2$; $I_{sp}=1500$~s & Hipótese conservadora~\cite{r29}\\
P14 & Geometria de eclipse & sombra cilíndrica; LTAN 06:00/12:00 & Modelo~\cite{r25,r26}\\
P15 & Falhas & $\lambda=1{,}97\times10^{-5}$~h$^{-1}$; $\varphi$ = 5--30\%; $N_{req}=1000$; 99\% & Taxa de~\cite{r36}; $\varphi$ e $N_{req}$ são hipóteses\\
P16 & Capacidade de lançamento & 20--25~t por voo ao SSO & Hipótese sem fonte primária\\
P17 & Piso de lançamento & $\Delta v$ 9,0--9,5~km/s; $I_{sp}$ e misturas da Seção~\ref{sec:lanc} & Hipótese~\cite{r30}\\
P18 & Monte Carlo & triangular, independente; $\rho_{rad}$: 4--22, moda 6 & Hipótese; variantes na Tabela~\ref{tab:mcvar}\\
P19 & WUE & 0,32--0,67 (base); 1--2; 5~L/kWh & Base e extremo de~\cite{r8}; 1--2 é hipótese dos autores, compatível em ordem de grandeza com~\cite{r22,r5}\\
P20 & Constantes do radiador e das perturbações & $\varepsilon=0{,}88$; $\eta_{fin}=0{,}94$; $T=328$~K; IR 237~W/m$^2$; albedo 0,30; $\alpha_s=0{,}2$; fator de iluminação 0,5; braço aerodinâmico 5~m; dipolo $10^{3}$~A\,m$^2$ & Hipóteses dos autores\\
P21 & Parâmetros de~\cite{r19} (Tabela~\ref{tab:cal19}) & manta 3,41~kg/m$^2$; bateria 3,75~kg/kWh; radiador 3,03~m$^2$/kW $\times$ 10~kg/m$^2$ & Tabela~1 de~\cite{r19}\\
P22 & Lançamento + reentrada & 51,7 (Falcon~9); 34,4~kg $\COeq$/kg (Starship) & \cite{r18}, v3\\
\end{longtable}

% ---------------------------------------------------------------------
\begin{thebibliography},{99}
\small

\bibitem{r1} ISO. \emph{ISO 14040:2006 --- Environmental management --- Life cycle assessment --- Principles and framework}. Geneva: ISO, 2006.
\bibitem{r2} ISO. \emph{ISO 14044:2006 --- Environmental management --- Life cycle assessment --- Requirements and guidelines} (inclui Amd 1:2017 e Amd 2:2020). Geneva: ISO.
\bibitem{r3} AIAA. \emph{AIAA S-120A-2015: Mass Properties Control for Space Systems}. Reston, VA, 2015.
\bibitem{r4} ECSS. \emph{ECSS-E-ST-10-02C Rev.1: Space engineering --- Verification}. Noordwijk, 2018.
\bibitem{r5} IEA. \emph{Energy and AI}. Paris, 2025.
\bibitem{r6} SHEHABI, A. et al. \emph{2024 United States Data Center Energy Usage Report}. LBNL, 2024. DOI 10.71468/P1WC7Q.
\bibitem{r7} SMITH, S. J. et al. \emph{United States Data Center Energy Usage Report: 2025 Update}. LBNL, 2026. DOI 10.71468/P1RP4F.
\bibitem{r8} GRÖGER, J.; BEHRENS, F.; GAILHOFER, P.; HILBERT, I. \emph{Environmental Impacts of Artificial Intelligence}. Öko-Institut, 2025.
\bibitem{r9} NASA. \emph{Total Solar Irradiance}. NASA Goddard.
\bibitem{r10} SPACEX. \emph{Falcon 9: Vehicle overview}.
\bibitem{r11} SPACEX. \emph{Falcon Heavy: Vehicle overview}.
\bibitem{r12} SPACEX. \emph{Starship: Design and vehicle overview}.
\bibitem{r13} MALONEY, C. et al. The climate and ozone impacts of black carbon emissions from global rocket launches. \emph{J. Geophys. Res. Atmos.}, v. 127, n. 12, e2021JD036373, 2022.
\bibitem{r14} RYAN, R. G. et al. Impact of rocket launch and space debris air pollutant emissions on stratospheric ozone and global climate. \emph{Earth's Future}, v. 10, n. 6, e2021EF002612, 2022.
\bibitem{r15} BROWN, T. F. M. et al. Worldwide rocket launch emissions 2019: An inventory for use in global models. \emph{Earth and Space Science}, v. 11, n. 10, e2024EA003668, 2024.
\bibitem{r16} FERREIRA, J. et al. Potential ozone depletion from satellite demise during atmospheric reentry in the era of mega-constellations. \emph{Geophys. Res. Lett.}, v. 51, n. 11, e2024GL109280, 2024.
\bibitem{r17} AL AHMAD, M.; MEMON, Q.; PECHT, M. Reliability and risk in space-based data centers: A lifecycle assessment of orbital cloud infrastructure. \emph{Applied Sciences}, v. 16, n. 11, 5247, 2026.
\bibitem{r18} OHS, R.; STOCK, G. F.; SCHMIDT, A.; FRAIRE, J. A.; HERMANNS, H. Dirty Bits in Low-Earth Orbit: The Carbon Footprint of Launching Computers. \emph{ACM SIGENERGY Energy Informatics Review}, v. 5, n. 2, jul. 2025. DOI 10.1145/3757892.3757896. arXiv:2508.06250v3, 18 fev. 2026. Disponível em: \url{https://arxiv.org/pdf/2508.06250}.
\bibitem{r19} OHS, R.; STOCK, G. F.; SCHMIDT, A.; FRAIRE, J. A.; OTT, J.; HERMANNS, H. Dark Clouds Rising in Low-Earth Orbit: On Environmental Limits to Massive Orbital AI. In: \emph{LEO-NET '26: 4th Workshop on LEO Networking and Communication (ACM SIGCOMM 2026)}, Denver, CO, 17--21 ago. 2026. ACM, 6 p. DOI 10.1145/3789240.3827598. Disponível em: \url{https://dl.acm.org/doi/10.1145/3789240.3827598}.
\bibitem{r20} NOVASPACE. \emph{Space Data Centers: The Next Frontier for AI and Cloud Computing}. White paper, maio 2026.
\bibitem{r21} SANTOS, R. et al. Data Centers, Minerais Críticos, Energia e Geopolítica: as bases da Inteligência Artificial. \emph{Sociedade \& Natureza}, v. 37, e72849, 2025.
\bibitem{r22} JORGE, M. do A. Especialistas alertam para alto consumo de água por \emph{data centers} em meio a incentivos do governo para atração de projetos. \emph{Jornal da Unesp}, 20 mar. 2026. Disponível em: \url{https://jornal.unesp.br/2026/03/20/especialistas-alertam-para-alto-consumo-de-agua-por-data-centers-em-meio-a-incentivos-do-governo-para-atracao-de-projetos/}.
\bibitem{r23} PETERS, A. The environmental cost of putting data centers in space. \emph{Fast Company}, 13 maio 2026. Disponível em: \url{https://www.fastcompany.com/91540754/the-environmental-cost-of-putting-data-centers-in-space}.
\bibitem{r24} MOON, H. Musk and Bezos push orbital data centers as US groups warn of pollution and ozone risks. \emph{DongA Science}, 24 jul. 2026. Disponível em: \url{https://www.dongascience.com/en/news/79087}.
\bibitem{r25} VALLADO, D. A. \emph{Fundamentals of Astrodynamics and Applications}. 4. ed. Microcosm Press, 2013.
\bibitem{r26} WERTZ, J. R.; EVERETT, D. F.; PUSCHELL, J. J. (eds.). \emph{Space Mission Engineering: The New SMAD}. Microcosm Press, 2011.
\bibitem{r27} WIE, B. \emph{Space Vehicle Dynamics and Control}. 2. ed. AIAA, 2008.
\bibitem{r28} MONTENBRUCK, O.; GILL, E. \emph{Satellite Orbits}. Springer, 2000.
\bibitem{r29} PICONE, J. M. et al. NRLMSISE-00 empirical model of the atmosphere. \emph{J. Geophys. Res.}, v. 107, n. A12, 1468, 2002.
\bibitem{r30} SUTTON, G. P.; BIBLARZ, O. \emph{Rocket Propulsion Elements}. 9. ed. Wiley, 2017.
\bibitem{r31} MARKLEY, F. L.; CRASSIDIS, J. L. \emph{Fundamentals of Spacecraft Attitude Determination and Control}. Springer, 2014.
\bibitem{r32} FRANKLIN, G. F.; POWELL, J. D.; WORKMAN, M. L. \emph{Digital Control of Dynamic Systems}. 3. ed. Addison-Wesley, 1998.
\bibitem{r33} ERICKSON, R. W.; MAKSIMOVIĆ, D. \emph{Fundamentals of Power Electronics}. 3. ed. Springer, 2020.
\bibitem{r34} MIDDLEBROOK, R. D. Input filter considerations in design and application of switching regulators. \emph{IEEE IAS Annual Meeting}, 1976.
\bibitem{r35} HASTINGS, D.; GARRETT, H. \emph{Spacecraft--Environment Interactions}. Cambridge University Press, 1996.
\bibitem{r36} DUBEY, A. et al. The Llama 3 Herd of Models. arXiv:2407.21783, 2024.
\bibitem{r37} GUINÉE, J. B. et al. \emph{Handbook on Life Cycle Assessment}. Kluwer, 2002.
\bibitem{r38} HELLWEG, S.; MILÀ I CANALS, L. Emerging approaches, challenges and opportunities in life cycle assessment. \emph{Science}, v. 344, n. 6188, p. 1109--1113, 2014.
\bibitem{r39} BOULAY, A.-M. et al. The WULCA consensus characterization model for water scarcity footprints (AWARE). \emph{Int. J. Life Cycle Assess.}, v. 23, p. 368--378, 2018.
\bibitem{r40} MASANET, E. et al. Recalibrating global data center energy-use estimates. \emph{Science}, v. 367, n. 6481, p. 984--986, 2020.
\bibitem{r41} ISO/IEC 30134-2:2016. \emph{Data centres --- Key performance indicators --- Part 2: Power usage effectiveness (PUE)}.
\bibitem{r42} ISO/IEC 30134-9:2022. \emph{Data centres --- Key performance indicators --- Part 9: Water usage effectiveness (WUE)}.
\bibitem{r43} ISO 24113:2019. \emph{Space systems --- Space debris mitigation requirements}.
\bibitem{r44} ECSS-E-ST-10-04C. \emph{Space engineering --- Space environment}.
\bibitem{r45} ECSS-E-ST-31C. \emph{Space engineering --- Thermal control general requirements}.
\bibitem{r46} ECSS-E-ST-60-30C. \emph{Space engineering --- Satellite attitude and orbit control system (AOCS) requirements}.
\bibitem{r47} ECSS-E-ST-20-08C. \emph{Space engineering --- Photovoltaic assemblies and components}.
\bibitem{r48} TRIBBLE, A. C. \emph{The Space Environment}. Princeton University Press, 2003.
\bibitem{r49} LEFFERTS, E. J.; MARKLEY, F. L.; SHUSTER, M. D. Kalman filtering for spacecraft attitude estimation. \emph{J. Guidance, Control, and Dynamics}, v. 5, n. 5, p. 417--429, 1982.
\bibitem{r50} DARKSKY INTERNATIONAL; ENVIRONMENT AMERICA; PUBLIC EMPLOYEES FOR ENVIRONMENTAL RESPONSIBILITY. \emph{Petition to Initiate Programmatic Environmental Impact Statement on Orbital Data Center Applications}. Petição perante a Federal Communications Commission, Washington, D.C., 8 jul. 2026. Representantes: J. E. Hasselman e S. W. Stern (Earthjustice). Disponível em: \url{https://earthjustice.org/document/space-data-center-petition}.

\end{thebibliography}

\end{document}
